% ==================================================================
% Conscious Point Physics:
% The Hidden Variable Factor K_0(lambda):
% Single-CP Path Integral and the Non-Perturbative Nexus Programme
% CPP Companion SD#5 — Version 0 (Research Agenda)
%
% STATUS: PLACEHOLDER / RESEARCH AGENDA PAPER
% This paper records the open problem identified at the close of
% SD#4 and specifies the programme required to resolve it.
% It is a formal reminder and agenda document, not a completed
% derivation. Version 1 will contain the actual calculation.
%
% GitHub: CPP/series_foundations/cpp_sd5_K0_derivation_v0.tex
% ==================================================================

\documentclass[11pt,a4paper]{article}
\usepackage[margin=1in]{geometry}
\usepackage[utf8]{inputenc}
\usepackage[T1]{fontenc}
\usepackage{amsmath,amssymb,amsthm}
\usepackage{caption,float,booktabs,parskip,xcolor}
\usepackage[numbers,sort&compress]{natbib}
\usepackage{hyperref}
\hypersetup{colorlinks=true,linkcolor=blue,urlcolor=cyan,citecolor=red}

\newtheorem{theorem}{Theorem}[section]
\newtheorem{proposition}[theorem]{Proposition}
\newtheorem{conjecture}[theorem]{Conjecture}
\newtheorem{remark}[theorem]{Remark}
\newtheorem{openprob}[theorem]{Open Problem}

\newcommand{\lP}{l_P}
\newcommand{\phig}{\varphi}
\newcommand{\eps}{\varepsilon}
\newcommand{\Napp}{N_{\rm app}}
\newcommand{\Npart}{N_{\rm part}}
\newcommand{\ket}[1]{|#1\rangle}

% Prominent status box
\newcommand{\statusbox}[1]{%
  \begin{center}
  \fboxsep=8pt
  \fbox{\parbox{0.85\textwidth}{\centering\large\bfseries #1}}
  \end{center}
  \vspace{6pt}
}

\title{\textbf{Conscious Point Physics:\\
The Hidden Variable Factor $K_0(\lambda)$:\\
Single-CP Path Integral and the\\
Non-Perturbative Nexus Programme}\\[6pt]
{\large Companion SD\#5 --- Version 0}\\[4pt]
{\normalsize\color{red}\bfseries Research Agenda Paper ---
Not a Completed Derivation}}

\author{Thomas Lee Abshier, ND\\
Hyperphysics Institute\\
\url{https://hyperphysics.com}\\
\texttt{drthomas007@protonmail.com}}

\date{23 March 2026 (agenda recorded); Version~1 target: TBD}

\begin{document}
\maketitle

\statusbox{STATUS: OPEN PROBLEM AGENDA\\
This is a v0 placeholder. The derivation is not yet complete.\\
This document records exactly what must be done and why it matters.}

\begin{abstract}
SD\#4~\citep{cpp_sd4} completed the CPP superdeterminism series
with two exact theorems (macroscopic limit $K \to 0$; single-CP
limit $K = K_0 f_{H_4}$) and one conjecture (the interpolation
formula).  Two quantitative questions were left open:

\begin{enumerate}
\item \textbf{The hidden variable factor $K_0(\lambda)$:} What is
  the explicit functional form of $K_0(\lambda)$, the projection
  of the particle hidden variable onto the DI-bit exchange degree
  of freedom?

\item \textbf{Non-perturbative proof of the interpolation:}
  The conjecture $K \approx \eps \cdot K_0 \cdot f_{H_4}$ is
  supported by a saddle-point calculation but not yet proved
  to all orders in $\eps$.
\end{enumerate}

These two questions determine the \emph{amplitude} of the CPP
correction to Bell correlations.  The angular form is already exact
(SD\#2) and the ratio test $\delta E(36°)/\delta E(120°) \approx
-1.065$ does not depend on the amplitude (SD\#2 eq.~21).
Nevertheless, a complete theory must provide the amplitude
$A_5 = \phig^{-3}/(2\pi)$ (conjectured) from first principles.

This paper records the mathematical programme required to answer
both questions, identifies the specific techniques available from
existing CPP results, and states what new mathematics is needed.
It is a \emph{research agenda} for future work, not a completed
derivation.  Upon completion of the derivation, this document
becomes Version~1 with the full results.
\end{abstract}

\tableofcontents
\newpage

%=====================================================================
\section{Why This Open Problem Matters}
%=====================================================================

The CPP prediction for Bell test deviations is (SD\#4 eq.~9):
\begin{equation}
E_{\rm CPP}(\theta) = -\cos\theta
+ \eps\,\bar{K}_0\,f_{H_4}(\theta) + O(\eps^2),
\label{eq:E_CPP}
\end{equation}
where $\eps \approx 2\times10^{-26}$ is the Nexus participation
fraction (known) and $f_{H_4}(\theta) = A_5\cos5\theta +
A_3\cos3\theta + \ldots$ is the angular function from SD\#2 (form
known, amplitudes $A_5, A_3$ conjectured).

The factor $\bar{K}_0 = \langle K_0(\lambda)\cos\theta_\lambda
\rangle_\lambda$ is the expectation value of $K_0$ weighted by the
hidden variable distribution.  Until $K_0(\lambda)$ is derived:

\begin{itemize}
\item The \emph{amplitude} of the CPP signal is unknown (but bounded:
  $|\eps\bar K_0| \leq \eps \approx 2\times10^{-26}$).
\item The conjecture $A_5 = \phig^{-3}/(2\pi)$ cannot be proved
  or disproved.
\item The non-perturbative verification of the interpolation formula
  is unavailable.
\end{itemize}

What is already known without $K_0$:
\begin{itemize}
\item The angular structure $f_{H_4}(\theta)$ is exact.
\item The ratio $\delta E(36°)/\delta E(120°) \approx -1.065$
  is exact (depends only on $\phig$, not on $\bar K_0$).
\item The null prediction at $\theta = 90°$ is exact.
\item The scaling laws ($\propto N_q$, $\propto 1/T$) are exact.
\end{itemize}

The missing piece is the overall scale, not the shape.

%=====================================================================
\section{The Mathematical Problem}
%=====================================================================

\subsection{The single-CP path integral}

From SD\#4 Theorem~2, the single-CP Nexus path integral is:
\begin{equation}
Z_{\rm CPP}^{(1,1)}(\lambda,\theta)
= \int \mathcal{D}[\phi_P]\,
  e^{i(S_P[\phi_P] + S_A[-\phi_P] + S_{PA}[\phi_P,\theta])/\hbar},
\label{eq:Z_single}
\end{equation}
where the apparatus field is $\phi_A = -\phi_P$ (constrained),
and $S_{PA}$ is the DI-bit exchange coupling between the two CPs.

The $K_0(\lambda)$ factor is:
\begin{equation}
K_0(\lambda) = \frac{Z_{\rm CPP}^{(1,1)}(\lambda,\theta)}
                    {Z_{\rm QM}^{(1)}(\lambda) \cdot Z_{\rm QM}^{(1)}(\text{vac})}
               - 1 \bigg|_{\theta\text{-independent part}}.
\label{eq:K0_def}
\end{equation}
The $\theta$-independent part is the projection of the full
integrand onto the $\cos(5\theta)$ mode (the leading $f_{H_4}$ term);
the result is $K_0(\lambda)$.

\subsection{What makes this hard}

Three technical obstacles:

\textbf{1. The CPP action $S[\phi]$ at the single-Grid-Point level.}
The QM series (2040a~\citep{cpp2040a}) derives the Schrödinger
equation from $Z_{\rm QM}$ in the coarse-grained limit.  The
single-Grid-Point action $S[\phi_i]$ at lattice resolution is
not yet written in closed form.  It is known to involve the
SSV field equations (companions C1--C2~\citep{abshier_sr}) and
the DI-bit density/phase Lagrangian, but the explicit form for
a single eCP at a single Grid~Point requires extracting the
lattice-scale limit of those field equations.

\textbf{2. The exchange coupling $S_{PA}[\phi_P, \theta]$.}
This is the interaction term between the particle CP and the
apparatus CP.  Its angular dependence on $\theta$ is what
generates the $f_{H_4}$ structure.  The derivation requires
computing the DI-bit exchange matrix element between two
adjacent Grid~Points in the 600-cell --- a calculation in
the adjacency operator formalism developed for the QM series
(companion 2040b~\citep{cpp2040b}).

\textbf{3. The Gaussian integration over $\phi_P$.}
Once $S_P + S_A + S_{PA}$ is written explicitly, the integral
over $\phi_P$ is a Gaussian path integral.  Evaluating it
gives $K_0$ as a function of the saddle-point value of
$\phi_P$, which depends on $\lambda$ through the initial
conditions.

\subsection{Candidate approaches}

\textbf{Approach A: Direct lattice calculation.}
Compute $S_{PA}$ explicitly for two adjacent 600-cell Grid~Points
in relative orientation $\theta$.  Evaluate the path integral in
the Gaussian approximation.  This gives $K_0$ perturbatively in
the DI-bit exchange coupling constant.

\textbf{Approach B: Transfer matrix.}
Express the Nexus path integral as a transfer matrix product
along the 600-cell adjacency graph.  For $N_{\rm part} = N_{\rm
app} = 1$, this is a $2 \times 2$ matrix problem.  The exact
result for the single-CP limit follows directly.

\textbf{Approach C: Symmetry argument.}
The $H_4$ symmetry of the 600-cell constrains $K_0(\lambda)$ to
transform in a specific irrep.  The simplest $H_4$-covariant
function with zero mean over $\lambda$ is:
\begin{equation}
K_0(\lambda) \sim \phig^{-3} \sin(\theta_\lambda),
\label{eq:K0_conjecture_form}
\end{equation}
which would give $\bar K_0 = \phig^{-3}/2$ and confirm the
conjecture $A_5 = \phig^{-3}/(2\pi)$.  This can be established
if the relevant $H_4$ irrep can be identified from the
representation theory of the 600-cell.

Approach C is the most tractable and should be attempted first.

%=====================================================================
\section{The Non-Perturbative Interpolation Problem}
%=====================================================================

\subsection{What needs to be proved}

SD\#4 Conjecture~1 states:
\begin{equation}
K(\lambda,\theta_A,\theta_B)
= \eps \cdot K_0(\lambda) \cdot f_{H_4}(\theta_A-\theta_B-\theta_\lambda)
+ O(\eps^2).
\label{eq:conj}
\end{equation}
The saddle-point calculation establishes the leading term.
The non-perturbative proof requires showing that all higher-order
terms $c_n \eps^n$ are bounded and that the series converges.

\subsection{Candidate approach: transfer matrix}

For the general case $\Npart = 2$, $\Napp \gg 1$, the Nexus
constraint distributes over $\Napp + N_{\rm rest}$ Grid~Points.
The transfer matrix along the 600-cell graph is:
\begin{equation}
T_{ij} = e^{i S_{\rm hop}(i \to j) / \hbar},
\label{eq:transfer}
\end{equation}
where $S_{\rm hop}$ is the single-hop DI-bit exchange action.
The Nexus partition function is:
\begin{equation}
Z_{\rm CPP} = \text{Tr}\!\left(T^{\Ntot}\right)\big|_{\rm constrained},
\label{eq:Z_transfer}
\end{equation}
where the constraint projects onto the zero-total-DI-bit sector.

For $\Napp \gg \Npart$, the spectral decomposition of $T^{\Napp}$
approaches a uniform distribution over the adjacency eigenvalues
(by the Perron--Frobenius theorem on the 600-cell graph).  This
gives the leading correction $K \sim \Npart/\Napp$ and the
convergence of the geometric series in~\eqref{eq:conj}.

This is a rigorous programme; it requires the spectral gap of the
600-cell transfer matrix, which is related to the adjacency
eigenvalue gap (already known from the QM/EW series:
$\lambda_{\rm max} = 12$, $\lambda_{\rm next} = 1+\phig \approx
2.618$).

%=====================================================================
\section{Research Programme and Milestones}
%=====================================================================

The following tasks, in order of recommended priority:

\begin{center}
\begin{tabular}{@{}clp{6cm}c@{}}
\toprule
Step & Task & Method & Feeds into \\
\midrule
1 & Identify $H_4$ irrep for $K_0$ & Representation theory of $H_4$
  (Fulton--Harris \citep{fulton1991}) & $A_5$ conjecture \\[4pt]
2 & Compute $S_{PA}(i,j,\theta)$ & SSV field equations (C2)
  + 600-cell adjacency operator & $K_0$ explicit \\[4pt]
3 & Evaluate single-CP Gaussian & Transfer matrix on 2-site
  system & $K_0$ confirmed \\[4pt]
4 & Spectral gap of 600-cell $T$ & Use $\lambda = \{12, 1+\phig,
  \phig-1, \ldots\}$ & Convergence proof \\[4pt]
5 & Non-perturbative proof of conj. & Transfer matrix + Perron--Frobenius
  & SD\#4 Conj.\ proved \\[4pt]
6 & Write SD\#5 v1 & Assemble steps 1--5 & Journal submission \\
\bottomrule
\end{tabular}
\end{center}

\subsection{What can be done now}

Step~1 (symmetry argument) requires only the representation theory
of $H_4$, which is in Humphreys~\citep{humphreys1990} and
Fulton--Harris~\citep{fulton1991}.  This is pure mathematics
requiring no new CPP calculations.

Step~4 (spectral gap) requires only the 600-cell eigenvalues
already computed in the QM and EW series.  The spectral gap is
$\lambda_{\rm max} - \lambda_{\rm next} = 12 - (1+\phig) =
11 - \phig \approx 9.382$.  This large gap implies rapid
convergence of the transfer matrix and supports the conjecture.

Steps~2 and~3 require the single-Grid-Point CPP action, which is
the main new mathematical input needed.

%=====================================================================
\section{Connection to the $A_5 = \phig^{-3}/(2\pi)$ Conjecture}
%=====================================================================

If the $H_4$ symmetry argument (Step~1) establishes
$K_0(\lambda) = \phig^{-3} \sin(\theta_\lambda)$, then:
\begin{equation}
\bar K_0
= \left\langle K_0(\lambda)\cos\theta_\lambda \right\rangle_\lambda
= \phig^{-3} \langle\sin\theta_\lambda\cos\theta_\lambda\rangle_\lambda
= \frac{\phig^{-3}}{2},
\label{eq:K0bar_if}
\end{equation}
using $\langle\sin\theta\cos\theta\rangle = 1/2$ for a uniform
distribution over $\theta_\lambda$.

The Fourier coefficient would then be:
\begin{equation}
A_5 = \bar K_0 / \pi = \frac{\phig^{-3}}{2\pi} \approx 0.0376,
\label{eq:A5_from_K0}
\end{equation}
confirming the SD\#2 conjecture $A_5 \approx \phig^{-3}/(2\pi)$
from first principles.

This would be a striking unification: the amplitude of the Bell
deviation is the same geometric factor $\phig^{-3} = V_{\rm
subgraph}/V_{\rm 600-cell}$ that governs the electroweak boson
masses (EW \#2~\citep{cpp_ew2}).  The EW masses and the quantum
foundations correction would be two manifestations of the same
600-cell geometric dilution.

This potential unification is the deepest motivation for completing
SD\#5.

%=====================================================================
\section{Open Problems (Formal Registry)}
%=====================================================================

\begin{openprob}[Explicit $K_0(\lambda)$ from $H_4$ symmetry]
\label{op:K0_symmetry}
Identify the $H_4$ irrep in which $K_0(\lambda)$ transforms.
Derive the explicit functional form from the representation
theory constraint.  Show that the conjectured form
$K_0(\lambda) \sim \phig^{-3}\sin\theta_\lambda$ is the unique
$H_4$-covariant, zero-mean, normalised function consistent with
the single-CP path integral structure.
\end{openprob}

\begin{openprob}[Single-CP CPP action $S[\phi_i]$]
\label{op:single_action}
Write the CPP action for a single eCP at a single Grid~Point in
the 600-cell lattice explicitly.  This requires taking the
$\Delta s \to \lP$ limit of the coarse-grained QM action
(companion 2040a~\citep{cpp2040a}) and expressing it in terms of
the DI-bit density $\rho_i$ and phase $\varphi_i$ at one site.
\end{openprob}

\begin{openprob}[DI-bit exchange matrix element $S_{PA}$]
\label{op:SPA}
Compute the DI-bit exchange coupling $S_{PA}[\phi_P,\theta]$
between two adjacent 600-cell Grid~Points in relative orientation
$\theta$.  This gives the angular dependence of the single-CP
path integral and confirms the $H_4$ structure from first
principles.
\end{openprob}

\begin{openprob}[Non-perturbative proof via transfer matrix]
\label{op:nonpert}
Apply the transfer matrix method on the 600-cell graph to prove
SD\#4 Conjecture~1 non-perturbatively.  Use the spectral gap
$\lambda_{\rm max} - \lambda_{\rm next} = 11 - \phig \approx 9.38$
to establish exponential convergence of the geometric series
$\sum_n c_n \eps^n$ in~\eqref{eq:conj}.
\end{openprob}

%=====================================================================
\section{Version History and Target}
%=====================================================================

\begin{center}
\begin{tabular}{@{}lll@{}}
\toprule
Version & Status & Content \\
\midrule
v0 (this document) & \textbf{Agenda only} & Problem statement,
  programme, formal open problems \\
v1 (target: TBD) & Full paper & Steps 1--5 completed;
  $K_0$ derived; $A_5$ confirmed or revised \\
\bottomrule
\end{tabular}
\end{center}

This v0 document is deliberately filed in the GitHub repository
so that the open problem is formally recorded, the mathematical
programme is specified, and future contributors can take it up
without repeating the analysis in SD\#1--4.

%=====================================================================
\section{Conclusion}
%=====================================================================

The CPP superdeterminism programme (SD\#1--4) is complete as a
qualitative theory.  The angular prediction, the scaling laws,
and the ratio test are exact.  What remains open is the
quantitative amplitude $\bar K_0$, which determines the overall
scale of the Bell deviation.

The key conjecture — $K_0(\lambda) \sim \phig^{-3}\sin\theta_\lambda$
and $A_5 = \phig^{-3}/(2\pi)$ — is accessible via the $H_4$
representation theory of the 600-cell without new physical input.
If confirmed, it would unify the electroweak mass sector (where
$\phig^{-3}$ is the geometric dilution of the subgraph volume) and
the quantum foundations sector (where $\phig^{-3}$ sets the Bell
deviation amplitude) under a single 600-cell lattice constant.

That unification is the prize that SD\#5 v1 should deliver.

\bibliographystyle{unsrtnat}
\begin{thebibliography}{12}

\bibitem{cpp_sd1}
T.~L. Abshier,
``The Nexus as superdeterministic mechanism,''
CPP SD\#1 v1, Hyperphysics Institute (2026).

\bibitem{cpp_sd2}
T.~L. Abshier,
``The 600-cell angular structure of Bell deviations,''
CPP SD\#2 v1, Hyperphysics Institute (2026).

\bibitem{cpp_sd3}
T.~L. Abshier,
``The CPP apparatus model,''
CPP SD\#3 v1, Hyperphysics Institute (2026).

\bibitem{cpp_sd4}
T.~L. Abshier,
``The Nexus correlation function,''
CPP SD\#4 v1, Hyperphysics Institute (2026).

\bibitem{cpp2040a}
T.~L. Abshier and Grok,
``Schrödinger equation from the 600-cell lattice,''
CPP-2040a v3.1, Hyperphysics Institute (2026).

\bibitem{cpp2040b}
T.~L. Abshier and Grok,
``Quantum statistics and the exchange interaction,''
CPP-2040b v3.1, Hyperphysics Institute (2026).

\bibitem{cpp_ew2}
T.~L. Abshier and Grok,
``The W$^0$ bracelet and W$^\pm$ boson,''
CPP EW \#2 v3.1, Hyperphysics Institute (2026).

\bibitem{hossenfelder_palmer}
S.~Hossenfelder and T.~Palmer,
``Rethinking superdeterminism,''
\textit{Front.\ Phys.} \textbf{8}, 139 (2020).

\bibitem{fulton1991}
W.~Fulton and J.~Harris,
\textit{Representation Theory: A First Course},
Springer (1991).

\bibitem{humphreys1990}
J.~E. Humphreys,
\textit{Reflection Groups and Coxeter Groups},
Cambridge University Press (1990).

\bibitem{brouwer1989}
A.~E. Brouwer, A.~M. Cohen, and A.~Neumaier,
\textit{Distance-Regular Graphs}, Springer (1989).

\bibitem{abshier_sr}
T.~L. Abshier,
``Mechanistic derivation of relativistic effects via Space Stress
Vector,'' viXra:2511.0062 (2025).

\end{thebibliography}

\end{document}
